%
% einleitung.tex -- Intro
%
% (c) 2020 Prof Dr Andreas Müller, Hochschule Rapperswil
%
% !TEX root = ../../buch.tex
% !TEX encoding = UTF-8
%
\section{Einleitung\label{qa:section:einleitung}}
\kopfrechts{Einleitung}
Mit Hilfe von imaginären Zahlen lassen sich viele Probleme der 
Mathematik elegant lösen. Dabei wird eine neue Regel eingeführt, 
welche die imaginäre Einheit $i$ mit der Eigenschaft $i^2 = -1$ 
definiert. Diese Einheit lässt sich, in Kombination mit reellen 
Zahlen, als eine neue Dimension interpretieren. Diese imaginäre 
Dimension steht rechtwinklig zum reelen Zahlenstrahl, sodass aus 
einer Linie eine Ebene entsteht. Die Ebene wird auch als komplexe 
Ebene bezeichnet. Zahlen, welche sowohl einen Imaginär- als auch 
einen Realteil beinhalten, gehören zur Menge der komplexen Zahlen 
$\mathbb{C}$. Sie bilden die Grundlage für wichtige praktische 
Anwendungen.

Doch was wäre, wenn weitere Dimensionen eingeführt werden könnten? 
Die Quaternionen, welche von William Rowan Hamilton im Jahr 1843 
entdeckt wurden, erweitern die reelen Zahlen um drei weitere 
Dimensionen $i,j,k$. Dafür wurde eine eigene Zahlenmenge $\mathbb{H}
$ geschaffen. Der Buchstabe $\mathbb{H}$ wurde zu Ehren von 
Hamilton gewählt.

Auf dem ersten Blick scheinen diese komplizierten Zahlen keiner 
praktischen Anwendung zuzuordnen. Jedoch spielen sie im heutigen 
digitalen Zeitalter eine grosse Rolle: In der Satellitennavigation, 
Raumfahrt, Quantenmechanik oder Gamingindustrie sind Quaternionen 
kaum mehr wegzudenken. Denn: Diese Bereiche benötigen Berechnungen 
von Orientierungen und Drehungen im 3D-Raum, wofür die Quaternionen 
bestens geeignet sind.


